\documentclass[10pt,a4paper]{amsart}
\usepackage[margin=19mm]{geometry}
\usepackage{amsmath,amssymb,mathtools}
\usepackage[scr=rsfs,cal=euler]{mathalfa}
\usepackage{xfrac}
\usepackage[unicode,hidelinks,bookmarks=false]{hyperref}
\newcommand{\ie}{i.e.\ }
\newcommand{\cf}{cf.\ }
\newcommand{\resp}{respectively\ }
\newcommand{\Subsection}{\subsection}
\providecommand{\nobreakdash}{\nobreak\hbox{-}}
\setlength{\emergencystretch}{2em}
\multlinegap0pt
\usepackage{bm}
\usepackage{amssymb}
\usepackage{booktabs}
\usepackage{multirow}
\usepackage{float}
\usepackage{xcolor}
\newtheorem{lemma}{Lemma}
\newtheorem{theorem}{Theorem}
\title{Fast Jacobi Spectral Methods and Closure Approximations for the Homogeneous FENE Model of Complex Fluids}
\author{Runkai Feng, Jie Shen, Haijun Yu}
\date{Source: arXiv:2602.07752v1 (CC BY 4.0)}
\begin{document}
\maketitle

\noindent\emph{Source and licence.} Fast Jacobi Spectral Methods and Closure Approximations for the Homogeneous FENE Model of Complex Fluids. Version arXiv:2602.07752v1 (CC BY 4.0), distributed by arXiv under the Creative Commons Attribution 4.0 International licence, http://creativecommons.org/licenses/by/4.0/, as recorded in the arXiv licence metadata for this exact version and shown on the abstract page https://arxiv.org/abs/2602.07752v1. Full licence notice: \texttt{licenses/arxiv-2602.07752-CC-BY-4.0.txt}.\par\smallskip

\noindent\textit{Editorial scope.} Counted unit: the article's theorem thm:stability (source0.tex lines 233--248). The article's complete original proof of it is delivered with it (source0.tex lines 250--283, one \texttt{proof} environment of 34 lines). No mathematics is written, repaired or re-derived: the statement and the proof are the article's own lines, in the article's own notation, and no retained line is substituted or re-flowed.  Retained uncounted as the source-local context the proof needs: lines 186--193; lines 201--215.  Nothing was omitted from any retained span, so the package carries no omission record.  No retained line contains a citation, so the wrapper carries no bibliography and no bibliography entry.  No retained line contains a figure environment, an image-inclusion command or drawing code, so no image is delivered and the package declares no asset dependency.  Editorial formatting only, in addition: an AMS \texttt{amsart} preamble that restates the article's own declarations and macros used by the retained lines, namely the environments \texttt{lemma}, \texttt{theorem} on separate counters, together with the standard packages those lines need.  The retained environments print in this document as Lemma 1, Theorem 1 in order of appearance, whereas the article prints the same items with the numbering of its own full document.  The complete native source of the article as distributed in the arXiv e-print for this exact version is bundled unchanged as \texttt{sources/194210/source0.tex}.
\begin{equation} 
\label{eq:BDF2}
\begin{aligned}
    & \frac{1}{2\Delta t}\left(3h^{n+1}-4h^{n}+h^{n-1}, g\right)_s + \frac{1}{\text{De}}\left(\nabla h^{n+1},\nabla g\right)_s \\
    &+\frac{b-2s}{\text{De}}\left(h^{n+1},\frac{\bm{q}\cdot \nabla g}{1-|\bm{q}|^2}\right)_s 
    = (2h^{n}-h^{n-1}, (\bm{K}\cdot\bm{q})\cdot\nabla g)_s.
\end{aligned}
\end{equation}

\begin{lemma}
\label{lem:Shen_identity}
For any $h \in H_s^1(\Omega)$ and $s > 1$, the following identity holds:
\begin{equation} \label{eq:identity}
    \left(h, \frac{\bm{q}\cdot \nabla h}{1-|\bm{q}|^2}\right)_s = (s-1)\|\bm{q}h\|_{s-2}^2 - \frac{3}{2}\|h\|_{s-1}^2.
\end{equation}
Furthermore, by analyzing the quadratic form associated with the weighted norms, there exists a threshold $\gamma(s)$ such that for any $\gamma > \gamma(s)$, the following inequality holds:
\begin{equation} \label{eq:positivity}
    \gamma \|h\|_s^2 + \left(h, \frac{\bm{q}\cdot \nabla h}{1-|\bm{q}|^2}\right)_s \ge \beta(\gamma)\|h\|_{s-2}^2,
\end{equation}
where 
\begin{equation}
    \gamma(s) = \frac{3}{4}\left(1 + \frac{s-1}{3} + \frac{3}{4(s-1)}\right), \quad \beta(\gamma) = (s-1)\left(1 - \frac{\gamma(s)}{\gamma}\right).
\end{equation}
\end{lemma}

\begin{theorem}[Stability]
\label{thm:stability}
Let $h^0 \in L^2_s(\Omega)$ be the initial data. Assume $1 < s \le b/2$. 
\begin{enumerate}
    \item \textbf{Case I ($s = b/2$):} The scheme \eqref{eq:BDF2} is unconditionally stable.
    \item \textbf{Case II ($s < b/2$):} The scheme is stable provided that the time step satisfies $\Delta t < \frac{3\mathrm{De}}{2(b-2s)\gamma(s)}$.
\end{enumerate}
In both cases, the solution satisfies the following energy estimate for any $N \ge 1$:
\begin{equation}
\begin{aligned}
    \|h^N\|_s^2 &+ \|2h^N - h^{N-1}\|_s^2 + \sum_{n=1}^{N-1} \|h^{n+1} - 2h^n + h^{n-1}\|_s^2 \\
    &+ \frac{2\Delta t}{\text{De}} \sum_{n=1}^{N-1} \|\nabla h^{n+1}\|_s^2 + 2\Delta t \sum_{n=1}^{N-1} \mathcal{P}(h^{n+1}) \le C(T, h^0, h^1),
\end{aligned}
\end{equation}
where $\mathcal{P}(h) = \beta(\gamma)\|h\|_{s-2}^2 \ge 0$ is the dissipation contribution from the FENE potential (with $\beta(\gamma)=0$ when $s=b/2$).
\end{theorem}

\begin{proof}
We choose the test function $g = h^{n+1}$ in \eqref{eq:BDF2}. 
Using the identity $2(3a-4b+c, a) = |a|^2 + |2a-b|^2 - |b|^2 - |2b-c|^2 + |a-2b+c|^2$, the first term becomes:
\begin{equation}
    \frac{1}{2\Delta t}\left(3h^{n+1}-4h^{n}+h^{n-1}, h^{n+1}\right)_s = \frac{1}{4\Delta t}\left( \mathcal{E}^{n+1} - \mathcal{E}^n + \|h^{n+1}-2h^n+h^{n-1}\|_s^2 \right),
\end{equation}
where $\mathcal{E}^{n+1} = \|h^{n+1}\|_s^2 + \|2h^{n+1}-h^n\|_s^2$.

For the convection term on the right-hand side, we apply the Cauchy-Schwarz inequality, Young's inequality, and the bound $\|\boldsymbol{K}\cdot\boldsymbol{q}\|_{L^\infty}^2 \le |\boldsymbol{K}|_\infty^2$:
\begin{equation}
\begin{aligned}
    \left( (2h^n - h^{n-1}) (\bm{K}\cdot \bm{q}), \nabla h^{n+1} \right)_s 
    &\le |\boldsymbol{K}|_\infty \|2h^n - h^{n-1}\|_s \|\nabla h^{n+1}\|_s \\
    &\le \frac{1}{2\text{De}} \|\nabla h^{n+1}\|_s^2 + \frac{\text{De}|\boldsymbol{K}|_\infty^2}{2} \|2h^n - h^{n-1}\|_s^2.
\end{aligned}
\end{equation}

Now we handle the FENE potential term.
\textbf{Case I ($s=b/2$):} The term involving $\frac{b-2s}{\text{De}}$ vanishes. We simply keep $\frac{3}{4\Delta t}\|h^{n+1}\|_s^2$ on the LHS.
\textbf{Case II ($s<b/2$):} We borrow a portion of the mass term from the time discretization to stabilize the potential term. Specifically, we set $\frac{3}{4\Delta t} = \frac{1}{4\Delta t} + \frac{b-2s}{\text{De}}\gamma$ (implying $\gamma = \frac{\text{De}}{2\Delta t(b-2s)}$). The stability condition $\Delta t < \frac{3\text{De}}{2(b-2s)\gamma(s)}$ ensures that $\gamma > \gamma(s)$. By Lemma \ref{lem:Shen_identity}, we have:
\begin{equation}
    \frac{b-2s}{\text{De}}\gamma \|h^{n+1}\|_s^2 + \frac{b-2s}{\text{De}}\left(h^{n+1}, \frac{\bm{q}\cdot \nabla h^{n+1}}{1-|\bm{q}|^2}\right)_s \ge  \frac{b-2s}{\text{De}}\beta(\gamma)\|h^{n+1}\|_{s-2}^2.
\end{equation}

Combining these estimates, we arrive at:
\begin{equation}
\begin{aligned}
    \frac{1}{4\Delta t} &(\mathcal{E}^{n+1} - \mathcal{E}^n) + \frac{1}{4\Delta t}\|h^{n+1}-2h^n+h^{n-1}\|_s^2 \\
    &+ \frac{1}{2\text{De}}\|\nabla h^{n+1}\|_s^2 +  \frac{b-2s}{\text{De}}\beta(\gamma)\|h^{n+1}\|_{s-2}^2 \\
    &\le \frac{\text{De}|\boldsymbol{K}|_\infty^2}{2} \|2h^n - h^{n-1}\|_s^2.
\end{aligned}
\end{equation}
Multiplying by $4\Delta t$, summing from $n=1$ to $N-1$, and noting that $\|2h^n - h^{n-1}\|_s^2 \le \mathcal{E}^n$, we can apply the discrete Gronwall inequality. Since the term $\beta(\gamma)\|h^{n+1}\|_{s-2}^2$ is non-negative, it is retained on the left-hand side to improve the stability bound. This completes the proof.
\end{proof}

\end{document}
